\section{Conclusion}\label{sec:conclusion}
    
    We have demonstrated that carefully crafted \textbf{data synthesis} can reshape the distribution of factual language corpora in a way that \emph{unlocks} grokking-based generalization. Even a moderately sized GPT-2 model can achieve substantial gains in \emph{multi-hop reasoning} by leveraging the late-phase formation of internal circuits -- outperforming more powerful models that do not receive synthetic data augmentation. Moreover, our empirical results indicate that factuality is not significantly compromised; on average, the model’s answers become more accurate when given a well-balanced mixture of real and synthesized facts. The main message of this work is that \textbf{boosting the inferred-to-atomic ratio $\phi_r$ via synthetic data remains the most direct route to emergent reasoning circuits}.
    
    Nevertheless, our approach also highlights several \textbf{limitations}. Full generalization across all relations requires each relation’s atomic facts to be sufficiently augmented, which can be challenging for \emph{rare} or \emph{low-frequency} relations. In many real-world corpora, knowledge graphs are not only sparse but also \emph{disconnected} or partially \emph{non-injective}, limiting the number of multi-hop paths the model can learn from. 
    
    Finally, \textbf{natural language challenges} persist in practical contexts. Real-world text often contains ambiguous references, unevenly distributed relations, and disjoint subgraphs, making high-quality data augmentation non-trivial. Nonetheless, our work illustrates the promise of \emph{implicit reasoning} -- once the ratio of inferred facts surpasses a critical threshold, the model internalizes robust logic circuits that can tackle complex multi-hop queries. We hope these findings encourage further research on efficient synthesis methods, broader domain applications, and deeper analysis of the mechanics behind grokking in large language models.
     